% ID: SEL-D0026-OA-E04
% ESTADO: cotejada. Segunda pasada visual completa y borrador compilado revisados.
% Fuente: Selectividad/matematicas_aplicadas/suplemente_septiembre_EXAMEN 4-ANDALUCÍA.pdf
% SHA-256: c41ea8ffe46f2de3a4d7f82e43def97322e29653182ad3df54a57e2ead2beb19
% Página: 1; opción A; ejercicio original 4.

Para estudiar el número de personas que van al cine mensualmente en una ciudad, se ha seleccionado una muestra aleatoria de 10 meses y se ha registrado el número de entradas al cine vendidas en cada mes. Los datos son los siguientes:
\[
682\quad553\quad555\quad666\quad657\quad649\quad522\quad568\quad700\quad552
\]
\begin{enumerate}[label=\alph*)]
\item \textbf{(2 puntos)} Suponiendo que el número de entradas vendidas mensualmente sigue una distribución Normal con desviación típica 50 entradas, calcule un intervalo de confianza, con un nivel del 95\%, para el número medio de personas que van al cine mensualmente en esa ciudad.
\item \textbf{(0.5 puntos)} \textquestiondown{}Cuál es el error máximo que se comete al estimar esta media con este intervalo?
\end{enumerate}
